\begin{pblock}{RTK GNSS field survey}
Ría de Villaviciosa, at the lowest spring tide of the month:
\term{361 fixed points at 1.3\,cm}. The implied datum offset between survey
and satellite is 52.81\,m — the local geoid undulation, an independent
confirmation that the two describe the same surface in different vertical
frames.

\vspace{3mm}
Each tide model is run through \emph{both} estimators, so every row compares
like with like, on the same 107 pixels. RMSE in metres:

\vspace{2mm}
\begin{tabular}{@{}l r r r l@{}}
\toprule
Tide model & {HSR} & {Step (DEA)} & {$\Delta$} & \\
\midrule
GOT4.10   & 0.159 & 0.135 &  0.024 & $p=0.007$ \\
EOT20     & 0.135 & 0.151 & -0.016 & n.s. \\
Ensemble  & 0.138 & 0.148 & -0.010 & n.s. \\
\bottomrule
\end{tabular}

\vspace{3mm}
\begin{pcallout}
\term{No configuration separates the estimators} under a paired bootstrap:
the tide model moves the numbers more than the estimator does. And
comparing GNSS \emph{points} to 10\,m pixel \emph{medians} adds
$\sigma\!\approx\!0.21$\,m of yardstick error — the method's own error is
$\approx$0.13--0.15\,m.
\end{pcallout}

\end{pblock}
